% % % % % % % % % % % % % % % % % % % % % % % % % % % % % % % % % % % % % % % % % % % % % % % % % % % % % % % % % % % % % % % %
% SAT_STYLE_MATH_MCQS_FR.tex
% 2_02.tex
%
% encoding: UTF-8
% EOL: LF
%
% format: LaTeX
% indent: spaces (2)
% width: 127
% % % % % % % % % % % % % % % % % % % % % % % % % % % % % % % % % % % % % % % % % % % % % % % % % % % % % % % % % % % % % % % %
Alice et Bernard jouent à pile ou face. Si pile, Alice reçoit $1$\,\currency\,de Bernard ;
sinon elle lui en doit un. Après (exactement) $10$ parties, Bernard a gagné $2$\,\currency. %
Le taux de succès d'Alice est de
%
\begin{enumerate}
  \item $60\%$
  \item $80\%$
  \item $90\%$
  \item $20\%$
  \item $40\%$.
\end{enumerate}
%
%
\begin{proof}%
Soient $N_A$ (respectivement $N_B$) les nombres de victoires d'Alice (respectivement Bernard). %
D'où 
%
\begin{equation}
  \begin{cases}
    N_A + N_B = 10, \\
    N_B - N_A = 2.
  \end{cases}
\end{equation}
%
Ainsi, %
%
\begin{equation}
  2N_A = N_A  + N_B - (N_B - N_A) = 10 -2 = 8, \textit{i.e., } N_A = 4.
\end{equation}
% 
Le taux de succès d'Alice est donc de 
%
\begin{equation}
  \frac{N_A}{N_A+ N_B} = \frac{4}{10} = 40\%.
\end{equation}
%
Réponse \textbf{E}.
\end{proof}